\section{Case Study: Existing Harnesses Through the Five Primitives}
\label{sec:casestudy}

We now test the position against harnesses in wide use. If it holds, each should decompose into the five primitives with nothing left over, and the differences between them should be locatable as decisions within a primitive.

\paragraph{Method.} We examined thirteen harnesses in three groups: command-line coding agents from model developers and the open-source community (Claude Code with its Agent SDK \citep{claudecode,agentsdkblog}, Codex CLI \citep{codexcli}, Gemini CLI \citep{geminicli}, Aider \citep{aider} and Cline \citep{cline}); research systems for software engineering (SWE-agent \citep{yang2024sweagent}, mini-SWE-agent \citep{minisweagent}, OpenHands \citep{wang2025openhands,wang2025openhandssdk} and Agentless \citep{xia2024agentless}); and general-purpose agent frameworks (LangGraph \citep{langgraph}, the OpenAI Agents SDK \citep{openaiagentssdk} and smolagents \citep{smolagents}, with the original AutoGPT \citep{autogpt} for historical comparison). For each, we read the primary sources (official documentation, source code and papers) and recorded how it implements each primitive and each production concern. Products in this space change quickly; the mapping describes them as documented in October 2026, and source-code observations refer to Codex CLI at commit \texttt{2f05c0a} and Gemini CLI at commit \texttt{42c5182}. Where we could not confirm a detail in a primary source, we omit it.

\subsection{Mapping}

Tables~\ref{tab:cli}, \ref{tab:research} and \ref{tab:frameworks} give the mapping. Every capability we found has a home in one primitive or in a production concern; we found no capability that required a sixth primitive.

\begin{table}[p]
\centering
\scriptsize
\setlength{\tabcolsep}{3pt}
\renewcommand{\arraystretch}{1.15}
\caption{Command-line coding agents, by primitive.}
\label{tab:cli}
\begin{tabularx}{\linewidth}{@{}>{\bfseries}p{1.95cm}YYYYY@{}}
\toprule
& Claude Code & Codex CLI & Gemini CLI & Aider & Cline \\
\midrule
\prim{Input} &
Interactive terminal; \texttt{-p} with piped stdin or a JSON stream; messages queued mid-turn; background output as events &
Interactive terminal; \texttt{codex exec} for scripts and CI, emitting JSONL events; SDK &
Interactive terminal with \texttt{@file} and \texttt{!}~shell; headless JSON output; ACP mode for editors &
Terminal chat with \texttt{code}, \texttt{ask}, \texttt{architect} and \texttt{help} modes &
Editor side panel; terminal output and linter or compiler errors \\
\addlinespace
\prim{Context} &
System prompt, tools, layered \texttt{CLAUDE.md} (or \texttt{AGENTS.md}), auto-memory, skill descriptions; clears old tool outputs then summarizes; deferred tool schemas loaded on demand &
Layered \texttt{AGENTS.md} from home to working directory, capped at 32\,KiB; compaction into a hand-off summary at a token limit &
Layered \texttt{GEMINI.md}; compression at half the window keeps the newest 30\% and summarizes the rest &
Repository map of signatures from tree-sitter, ranked by PageRank over a dependency graph, within a token budget; chat history summarized &
Project structure; automatic summarization near the limit, with rule-based truncation as fallback \\
\addlinespace
\prim{Model}\newline\prim{interface} &
Claude models via the Anthropic API or cloud providers; fallback model chain; streaming &
OpenAI Responses API, streamed; other providers through configurable model providers &
Gemini models via the Google GenAI SDK, streamed; optional local routing &
Many providers through LiteLLM; edits parsed from text, not tool calls &
Many providers; native tool-call schemas \\
\addlinespace
\prim{Output} &
About twenty tools: shell, read, exact-string edit, write, search, web, sub-agent, plan mode, notebook, language server; MCP tools &
Shell with resumable sessions, grammar-constrained \texttt{apply\_patch}, plan updates, image view, web search, multi-agent tools; MCP tools &
Shell, file read and write, exact-string replace, glob and grep, web search and fetch, to-do and plan tools; MCP tools &
Edit formats: whole file, search/replace, unified diff; architect mode pairs a planning model with an editing model &
About twenty tools: command, read, write, replace, patch, search, browser, web, MCP, sub-agents, completion \\
\addlinespace
\prim{Control flow} &
Loop until a reply has no tool calls; optional turn and budget limits; sub-agents with their own context; plan mode &
Loop over the Responses API; sub-agents on by default; plan mode; persistent goal &
Loop with a hard cap of 100 turns and loop detection; sub-agents exposed as tools; plan mode &
One edit per request, then lint and test, retrying failures up to three times &
Plan and Act modes, each with its own model; ends on \texttt{attempt\_completion}; sub-agents \\
\bottomrule
\end{tabularx}
\end{table}

\begin{table}[p]
\centering
\scriptsize
\setlength{\tabcolsep}{3pt}
\renewcommand{\arraystretch}{1.15}
\caption{Research systems for software engineering, by primitive.}
\label{tab:research}
\begin{tabularx}{\linewidth}{@{}>{\bfseries}p{1.95cm}YYYY@{}}
\toprule
& SWE-agent & mini-SWE-agent & OpenHands & Agentless \\
\midrule
\prim{Input} &
An issue or task statement, with optional demonstrations &
A task string rendered into templates &
User messages and environment observations in one event stream; issue triggers &
An issue and its repository \\
\addlinespace
\prim{Context} &
Templates for each step; history processors that keep the last five observations and collapse older ones &
Linear history: every step appends, nothing is processed &
Chronological event history; a summarizing condenser keeps the first events and a recent tail and summarizes the middle; \texttt{AGENTS.md} and skills loaded by trigger &
Hierarchical localization: repository tree, then file skeletons, then edit locations; embedding retrieval in parallel \\
\addlinespace
\prim{Model}\newline\prim{interface} &
LiteLLM; pluggable parsers for thought-and-action text or native tool calls &
LiteLLM and other gateways; no tool-call interface, so any model works &
LiteLLM; tool calls parsed from text for models without native support &
Plain completions; no tool calls \\
\addlinespace
\prim{Output} &
An agent-computer interface: file viewer of 100 lines, search with capped results, an edit command that rejects edits failing a linter &
Bash only, each command in a fresh subprocess &
Code actions: Python and shell in a sandbox; a browser; editing utilities &
Search/replace patches; forty candidate patches per issue \\
\addlinespace
\prim{Control flow} &
Open loop ending on \texttt{submit}; retries after malformed output; a cost budget that auto-submits; a meta-agent that picks the best of several attempts &
A \texttt{while} loop ending on a completion marker; step, cost and time limits &
A step loop; delegation to specialized agents; stuck detection &
A fixed pipeline: localize, repair, validate; a majority vote over normalized patches \\
\bottomrule
\end{tabularx}
\end{table}

\begin{table}[p]
\centering
\scriptsize
\setlength{\tabcolsep}{3pt}
\renewcommand{\arraystretch}{1.15}
\caption{Agent frameworks, by primitive, with AutoGPT (2023) for historical comparison.}
\label{tab:frameworks}
\begin{tabularx}{\linewidth}{@{}>{\bfseries}p{1.95cm}YYYY@{}}
\toprule
& LangGraph & OpenAI Agents SDK & smolagents & AutoGPT (2023) \\
\midrule
\prim{Input} &
Graph invocation with an initial state; resumption values after an interrupt &
Run input to a runner; session history; approvals to resume a paused run &
A task string; outputs of managed agents &
A goal set by the user; per-command authorization \\
\addlinespace
\prim{Context} &
A user-defined state schema, merged by per-key reducers; a message list with an ID-aware reducer; a long-term store &
Instructions, input and tool results only; local run context is explicitly \emph{not} sent to the model; handoff input filters &
Memory written to messages each step; managed agents described in the manager's system prompt; optional planning steps &
Recent history plus vector-memory retrieval, fitted to the token limit \\
\addlinespace
\prim{Model}\newline\prim{interface} &
Left to user code inside nodes &
Responses or Chat Completions models; custom providers; beta LiteLLM adapter &
Many model classes, including local inference and LiteLLM &
OpenAI chat API; JSON responses with thoughts and a command \\
\addlinespace
\prim{Output} &
Left to user code inside nodes; tool nodes &
Hosted tools, local shell and patch tools, function tools from signatures, agents as tools; approval before chosen tools &
Code actions (Python) or JSON tool calls; sandboxed executors &
Named commands with arguments (files, web, memory) \\
\addlinespace
\prim{Control flow} &
An explicit graph: nodes, conditional edges, commands that route or hand off to a parent graph, fan-out; super-step limit &
A runner loop: model call, then final output, handoff, or tool calls; handoffs as tools; turn limit &
A multi-step loop ending on \texttt{final\_answer} or a step limit; answer checks; managed agents &
An infinite loop with human authorization per command; a continuous mode with a limit \\
\bottomrule
\end{tabularx}
\end{table}


\subsection{Frameworks expose the primitives}

Agent frameworks differ from complete harnesses in \emph{which} primitives they implement and which they leave to the developer, and so they show the primitives directly. LangGraph describes itself as focused entirely on orchestration: it provides \prim{Control flow} as an explicit graph, and \prim{Context} as typed state merged by reducers, while model calls and tool execution happen in user-written nodes \citep{langgraph}. The OpenAI Agents SDK provides a fixed \prim{Control flow}, a runner loop that calls the model and then either finishes, hands off to another agent, or runs tool calls and loops, and draws an explicit line in \prim{Context}: local run state ``is not sent to the LLM'', which sees only the conversation \citep{openaiagentssdk}. smolagents frames agency itself as a matter of \prim{Control flow}, from a router to a multi-step loop to multiple agents, and exposes the main \prim{Output} decision as a choice of class: \texttt{CodeAgent} writes actions as Python, \texttt{ToolCallingAgent} as JSON tool calls, following \citet{wang2024codeact} \citep{smolagents}. AutoGPT, an early autonomous agent from 2023, already had all five: a goal as input, history plus retrieved vector memory as context, a JSON response of thoughts and a command, named commands as output, and an unbounded loop gated by human authorization \citep{autogpt}. Its safety model, a person approving each command unless a ``continuous mode'' was enabled, is the same decision that current harnesses express as approval modes.

\subsection{The differences are implementation decisions}

Reading the tables by row rather than by column shows each difference as a choice along an axis within one primitive.

\paragraph{\prim{Input}: where requests come from.} All of the coding agents accept both an interactive person and a non-interactive caller (a command-line flag, a JSON event stream or an SDK), so the same harness serves a developer and a continuous-integration job. Research systems take a task statement, usually an issue. These are different sources feeding the same primitive.

\paragraph{\prim{Context}: what the model is shown.} Three designs recur. The first is a \emph{project instruction file} merged from several directory levels: \texttt{CLAUDE.md}, \texttt{AGENTS.md} and \texttt{GEMINI.md} are the same decision under three names, and Claude Code and OpenHands read \texttt{AGENTS.md} as well. The second is a strategy for \emph{keeping history within the window}, and here the axis is clearest: mini-SWE-agent processes nothing; SWE-agent elides old observations, which its authors found better than keeping full history \citep{yang2024sweagent}; Claude Code, Codex CLI, Gemini CLI, OpenHands and Cline summarize, differing in when they trigger and what they keep verbatim. The third is \emph{what of the project to show}, from nothing (the agent explores through its tools) to Aider's ranked repository map to Agentless's staged localization. All are answers to one question: what should the model be shown on this call?

\paragraph{\prim{Model interface}: one provider or many, native or parsed.} Model developers' harnesses bind to their own API, while research and community harnesses route through a translation layer, most often LiteLLM, to reach many providers. A second axis is whether actions arrive as native tool calls or are parsed from text; mini-SWE-agent, Aider and Agentless parse text, and SWE-agent and OpenHands support both. Both axes trade portability against the use of provider-specific features.

\paragraph{\prim{Output}: how many tools, and in what format.} The axis runs from a single shell (mini-SWE-agent), through a small purpose-built interface (SWE-agent), to rich sets of around twenty tools (Claude Code, Cline, Gemini CLI) and to code as the action format (OpenHands, following \citet{wang2024codeact}). Edit formats are a sub-axis: exact string replacement (Claude Code, Gemini CLI, Cline), a grammar-constrained patch language (Codex CLI), and search/replace or unified diffs parsed from text (Aider, Agentless). Each harness also extends its tool set through the Model Context Protocol, which is a decision about how \prim{Output} discovers tools, not a new primitive.

\paragraph{\prim{Control flow}: who decides what happens next.} At one end, Agentless fixes the sequence of model calls in code; at the other, Claude Code, Codex CLI and mini-SWE-agent run until the model stops requesting actions. Between them lie explicit completion actions (\texttt{submit}, \texttt{attempt\_completion}), retry loops driven by tests (Aider), best-of-several attempts (SWE-agent's meta-agent), delegation to sub-agents with separate contexts, and modes that change what is allowed (plan modes, Aider's architect mode). Stopping rules vary in the same way: a hard turn cap (Gemini CLI), optional turn and budget limits (Claude Code's SDK), and cost or time limits (mini-SWE-agent, SWE-agent).

\subsection{Production concerns}

Table~\ref{tab:concerns-cli} shows the production concerns of the coding agents. The same pattern holds: each concern is answered inside a primitive, and the harnesses differ in their answers rather than in which concerns they have.

\begin{table}[t]
\centering
\scriptsize
\setlength{\tabcolsep}{3pt}
\renewcommand{\arraystretch}{1.15}
\caption{Production concerns in command-line coding agents.}
\label{tab:concerns-cli}
\begin{tabularx}{\linewidth}{@{}>{\bfseries}p{1.95cm}YYYYY@{}}
\toprule
& Claude Code & Codex CLI & Gemini CLI & Aider & Cline \\
\midrule
Persistence & Session transcripts, resume and fork; file checkpoints before each prompt with rewind & Session files, resume and fork; no built-in rewind (git recommended) & Auto-saved sessions, resume; checkpoints in a shadow git repository; rewind & Every change auto-committed to git; undo & Checkpoints in a shadow git repository after each tool use; restore files, task, or both \\
Safety \& security & Permission modes; deny/ask/allow rules; OS sandbox (Seatbelt, bubblewrap) with a network allow-list & Sandbox modes from read-only to full access, network off by default; approval policies; OS sandbox (Seatbelt, bubblewrap with seccomp) & Approval modes; a rule-based policy engine; sandboxes from Seatbelt to containers and gVisor; trusted folders & No sandbox; edits the working tree under version control & Approval for every edit and command, with per-category auto-approval; no sandbox \\
Observability & Hooks at about thirty events; OpenTelemetry & Hooks; OpenTelemetry & Hooks around agent, model and tool calls; OpenTelemetry & Commit history & Checkpoint history \\
Resilience & Configurable retries and timeouts; fallback models & Request and stream retries; idle timeout & Retries with backoff; model fallback & --- & --- \\
Performance & Prompt caching; read-only tools run concurrently & Prompt cache key; parallel tool calls & Token caching; parallel tool calls batched & Token-budgeted repository map & Context summarization \\
\bottomrule
\end{tabularx}
\end{table}

Two designs for safety stand out. \emph{Isolation} contains the damage an action can do: the model developers' harnesses use operating-system sandboxes, and the research systems run in containers. \emph{Reversibility} makes damage undoable: Aider commits every change, Cline and Gemini CLI keep a shadow git repository, and Claude Code checkpoints files. Most harnesses combine the two with approvals, which keep a person in the loop. These are different answers to the safety question, located in \prim{Output} and \prim{Control flow}.

\subsection{Two limit cases}

\paragraph{A harness without a loop.} Agentless is presented as an alternative to agents: a fixed pipeline of localization, repair and validation in which ``the LLM [does not] decide future actions'' \citep{xia2024agentless}. In our terms it is still a harness with all five primitives; its \prim{Control flow} is a fixed sequence written in code rather than an open loop, and its \prim{Context} does the heavy lifting through staged localization. Its competitive results at the time of publication show that the axis within \prim{Control flow} is a real design choice, not a ladder with agents at the top.

\paragraph{A harness with almost nothing.} mini-SWE-agent describes itself as a hundred-line agent: one tool (bash), a linear history with no processing, and no tool-call interface. Its authors report that it resolves over 74\% of SWE-bench Verified with a strong model \citep{minisweagent}. It shows the primitives at their smallest: each is present, and each is a single decision. That a harness this small is competitive, given a capable model, supports the position that the primitives, not any particular elaboration of them, are what make an agent.

\subsection{Convergence}

The three harnesses built by model developers, written independently, converged on nearly the same decisions: a layered instruction file in \prim{Context}, summarization to keep history in the window, a rich tool set extensible through MCP in \prim{Output}, sub-agents and plan modes in \prim{Control flow}, and, across the production concerns, operating-system sandboxes, rule-based permissions, hooks and OpenTelemetry. We read this convergence as evidence that the primitives pose the right questions: independent teams, answering the same five questions under similar constraints, arrive at similar answers.
